\section{A Value-Function View of RHO-LOSS}
\label{sec:connection}

\subsection{The F2SA value-function direction}

Define the lower-level value function
\begin{equation}
g^\star(\vomega):=\min_{\vtheta}g(\vomega,\vtheta)
=g(\vomega,\thetahat(\vomega)),
\qquad
\thetahat(\vomega)\in\argmin_{\vtheta}g(\vomega,\vtheta).
\end{equation}
The bilevel feasibility condition is
$g(\vomega,\vtheta)-g^\star(\vomega)\leq 0$. We therefore consider the
penalized value-function objective from the original formulation,
\begin{equation}
\mathcal{L}^{\alpha}(\vomega,\vtheta,\vtheta')
=f(\vtheta)+\alpha\bigl(g(\vomega,\vtheta)
-g(\vomega,\vtheta')\bigr),
\qquad \alpha>0,
\label{eq:f2sa-penalty}
\end{equation}
and track two solutions,
\begin{align}
\thetatilde(\vomega)
&\in\argmin_{\vtheta}
    \left[f(\vtheta)+\alpha g(\vomega,\vtheta)\right],
\label{eq:tilde-solution}\\
\thetahat(\vomega)
&\in\argmin_{\vtheta}g(\vomega,\vtheta).
\label{eq:hat-solution}
\end{align}
Under the regularity conditions needed for the envelope theorem, the gradient
of the minimized penalty is
\begin{equation}
\nabla_{\vomega}\mathcal{L}^{\alpha,\star}(\vomega)
=\alpha\left[
\nabla_{\vomega}g(\vomega,\thetatilde)
-\nabla_{\vomega}g(\vomega,\thetahat)
\right].
\label{eq:f2sa-gradient}
\end{equation}
For the linear weighting in Eq.~\eqref{eq:fg}, its $i$th coordinate is
\begin{equation}
\left[\nabla_{\vomega}\mathcal{L}^{\alpha,\star}(\vomega)\right]_i
=\frac{\alpha}{n}\left[
\ell(\thetatilde;z_i)-
\ell(\thetahat;z_i)
\right].
\label{eq:loss-difference-gradient}
\end{equation}
Consequently, the loss difference
\begin{equation}
s_i^{\mathrm{VF}}(\vomega)
:=\ell(\thetahat(\vomega);z_i)
-\ell(\thetatilde(\vomega);z_i)
\label{eq:vf-score}
\end{equation}
is proportional to the \emph{negative} penalized outer gradient. A selector
that keeps examples with large $s_i^{\mathrm{VF}}$ therefore follows a descent
direction, subject to inner-solution and finite-$\alpha$ errors. The same
$\alpha$ must appear in both the penalized inner objective and the outer
gradient. A fixed value is a finite-penalty approximation; a scheduled value
must be coordinated with the inner and outer step sizes. Score
standardization removes only the global scale of Eq.~\eqref{eq:f2sa-gradient},
not the dependence of $\thetatilde$ on $\alpha$.

\subsection{Ideal and practical RHO-LOSS}

RHO-LOSS begins with online batch selection. At training step $t$, a large
candidate batch $C_t$ is sampled, a score is computed for every example in
$C_t$, and only the top $b$ examples are used for one target-model update. Its
Bayesian derivation targets the example that would most reduce predictive
holdout loss. After replacing Bayesian posteriors with trained point-estimate
models, the idealized score can be written in our notation as
\begin{equation}
\rhoideal_i(\vomega)
=\ell(\thetahat(\vomega);z_i)
-\ell(\thetatilde(\vomega);z_i),
\label{eq:ideal-rho}
\end{equation}
where the second model incorporates both the holdout objective and the data
represented by $\vomega$. Comparing Eqs.~\eqref{eq:vf-score} and
\eqref{eq:ideal-rho} gives
\begin{equation}
\rhoideal(\vomega)
=s^{\mathrm{VF}}(\vomega)
=-\frac{n}{\alpha}
\nabla_{\vomega}\mathcal{L}^{\alpha,\star}(\vomega).
\label{eq:rho-vf-equivalence}
\end{equation}
This is an equality of directions at a fixed $\vomega$ under exact inner
solutions and aligned objective normalizations. It does not say that one
Top-$b$ RHO step solves Eq.~\eqref{eq:bilevel-selection}.

Practical RHO-LOSS makes an additional approximation. It trains an
irreducible-loss model once on $\Dval$,
\begin{equation}
\thetaIL\in\argmin_{\vtheta}f(\vtheta),
\end{equation}
freezes it, and uses
\begin{equation}
\rhopractical_i(t)
=\ell(\vtheta_t;z_i)-\ell(\thetaIL;z_i),
\label{eq:practical-rho}
\end{equation}
where $\vtheta_t$ is the current, generally nonconverged target model. This
removes the need to update the joint comparator as selected data change. It
also means Eq.~\eqref{eq:practical-rho} is not, in general, the exact gradient
in Eq.~\eqref{eq:rho-vf-equivalence}. The RHO-LOSS paper finds this frozen
approximation computationally useful and sometimes empirically preferable;
our framework treats that as a testable approximation rather than an
identity.

\subsection{The approximation ladder}

The shared formulation yields the following ladder:
\begin{equation}
\underbrace{\nabla h(\vomega)}_{\text{bilevel target}}
\;\longleftarrow\;
\underbrace{\nabla\mathcal{L}^{\alpha,\star}(\vomega)}_{\text{finite-penalty VF}}
\;\longleftrightarrow\;
\underbrace{-\rhoideal(\vomega)}_{\text{two exact inner models}}
\;\longrightarrow\;
\underbrace{-\rhopractical(t)}_{\text{frozen IL proxy}}.
\label{eq:approximation-ladder}
\end{equation}
The left arrow carries finite-penalty bias and relies on F2SA's regularity and
tracking conditions. The middle equivalence follows algebraically from
linear example weights. The right arrow introduces the frozen comparator and
target-trajectory approximations. Batch-local ranking and Top-$B$
discretization are further approximations applied after the score is formed.
